\subsection{General split construction and preservation}
\label{subsec:split-detail}
Let $U_1,\dots,U_r$ be the irredundant conjunctions of non-strict upper bounds
of a disjunctive invariant; the construction makes one copy per $U_k$. At
entry valuation $v$, the largest delay admitted by $U_k$ is
\[
 d_k(v)=\min_{(x,b)\in U_k}(b-v(x)),
 \qquad d_k(v)=+\infty\text{ when }U_k=\varnothing.
\]
The entry predicate for copy $k$ is
\[
 G_k(v)\equiv U_k(v)\ \land\ \bigwedge_{\ell=1}^{r} d_\ell(v)\le d_k(v).
\]
Thus a copy is admitted only if its invariant holds at entry and no other
copy permits a strictly longer stay. To encode the comparison without
introducing a new clock, the implementation chooses an upper bound
$y\le b_y$ in $U_\ell$ and requires it to be reached no later than every
bound $x\le b_x$ in $U_k$. This gives $v(x)-v(y)\le b_x-b_y$. The
disjunction over possible minimizing bounds is represented in DNF. When some
$U_\ell$ is empty, $d_\ell=+\infty$ and the comparison $d_\ell\le d_k$ holds only for an
empty $U_k$; the encoding then reduces to the constants true and false rather than to a choice of bound. An
incoming edge reset set $Z$ is handled by substituting zero for clocks in
$Z$ before conjoining $G_k$ with the edge guard: a difference constraint
with two reset clocks becomes a constant test, one reset clock gives a unary
comparison, and no reset clock leaves a diagonal constraint.

The copy invariants are generated from the outgoing transition guards. If
$\operatorname{ub}(c)$ extracts the non-strict unary upper bounds from a
conjunctive guard clause $c$, the coverage contract is
\[
\begin{gathered}
\forall t\in\operatorname{Out}(q),\ \forall c\in\operatorname{DNF}(g_t),\\
\exists U\in\mathcal U(q):\quad \operatorname{ub}(c)\Rightarrow U.
\end{gathered}
\]
Here $t$ has source $q$ and guard $g_t$; implication means that every
valuation satisfying all extracted upper bounds also satisfies every bound
in $U$. The Rocq lemma \texttt{disj\_inv\_cover} proves this coverage for
the deduplicated list $\mathcal U(q)$; it is an internal construction
invariant, not a precondition imposed on callers. At the endpoint of any taken
transition, the selected guard clause therefore entails a generated
upper-bound conjunction. Since clocks only increase during the stay, those
non-strict upper bounds also hold at each earlier valuation. This coverage
argument is the key completeness condition for both the verified split and
the all-copy comparison below.

All copies inherit the original location's accepting mark. No silent edge is
added, and the proof relates runs by projecting each copy back to its
original location. Thus repeated visits to an accepting location correspond
to repeated visits to accepting copies, preserving B\"uchi acceptance. In
Rocq, \texttt{split\_accepts} proves timed-word preservation when the input
invariant function is everywhere absent; \texttt{of\_tba} establishes that
precondition for the TBA export pipeline. The split invariant is absent or
one conjunction per location by \texttt{split\_inv\_conjunctive}; the full
export shape is established by \texttt{export\_guards\_conjunctive} and
\texttt{export\_invariants\_conjunctive}. These are the internal
representation theorems; the XML printer is not verified.

The optimized automaton has normalized conjunctive guards. Export first merges transitions that agree on source, target, event label, and resets, placing their guards into a disjunction. The location invariant is then represented by disjunctions of conjunctions of upper bounds. UPPAAL-oriented output requires each such invariant to be conjunctive because its location syntax admits a conjunction of clock conditions, so a disjunction must be encoded by a structural transformation~\cite{[BGD04]}. Retaining a disjunction directly would not meet the target representation.

Every incoming edge is replicated to each regular target copy and conjoined with its reset-adjusted entry predicate; outgoing edges are replicated from every copy. This replication is essential: copying only the location and assigning incoming edges to one arbitrarily chosen disjunct can truncate a stay and lose a timed word. The initial location uses a dedicated invariant-free copy because the first timed-word event has no incoming edge.

The representation cost follows directly from this construction. Let
$K=\mathsf{kp}(D)$ be the number of source-copy slots used by the
implementation and let $r(q)$ be the number of regular invariant copies of
location $q$. Before guard explosion, each original edge $t$ contributes
$K\,r(\mathit{tgt}(t))$ split transition records. After guard explosion,
the corresponding explicit-edge count is
\[
 E_{\mathrm{split}}=\sum_{t\in T}K\sum_{k<r(\mathit{tgt}(t))}
 \left|\operatorname{DNF}\!\left(g_t\land
   \operatorname{subst}_{Z_t}(G_k)\right)\right|.
\]
This count is taken before later normalization and pruning. It exposes two
sources of growth: the number of invariant copies at the target and the DNF
size of reset-adjusted entry guards. It is a structural count of the
construction, not a complexity bound for the entire exporter.

Language preservation does not require selecting a copy with the longest
admissible residence interval. We therefore mechanized a second policy,
called the \emph{all-copy split}: it retains the same conjunctive invariant
copies, invariant-free initial copy, and accepting marks, but sends every
incoming transition to every regular target copy with its original guard and
reset set. It has no reset-adjusted entry predicate. The copy coverage
property \texttt{disj\_inv\_cover} is established inside the Rocq
construction, not assumed from an external parser.

\begin{theorem}[All-copy split preservation]
Let $D$ be an input DTA with no pre-existing location invariants. The
all-copy construction computes its upper-bound disjunctions and proves
\texttt{disj\_inv\_cover} internally. Then the all-copy variant and $D$
accept the same timed words.
\end{theorem}
\begin{proof}
For soundness, project a run of the split automaton from each copy to its
original location. Each copied transition keeps its original label, guard,
and resets. The regular copies add their associated upper-bound conjuncts,
but the input automaton has no location invariants, so projection forgets
only this additional condition. The initial copy has no invariant, as required
by the input contract. Accepting marks are inherited, so infinitely many
accepting visits are preserved.

For completeness, fix an accepting run of $D$. At each noninitial location
occurrence, consider the outgoing transition used by that run and a DNF
clause $c$ of its guard that holds at the event ending the stay. By
\texttt{disj\_inv\_cover}, some copy invariant $U_k$ is implied by the
upper-bound part of $c$. Hence $U_k$ holds at that ending event. Each atom
of $U_k$ is a non-strict upper bound, so it also holds at every earlier
valuation in the stay: clocks only increase during delay, and resets occur
at its beginning. Thus copy $k$ satisfies its invariant throughout this
stay. The preceding transition is available to every regular target copy,
and every outgoing transition is available from each source copy. Choosing
copy $k$ on entry therefore lifts this occurrence of the run. The initial
occurrence uses the dedicated invariant-free copy. Making this choice at
each occurrence preserves the event sequence and resets; inherited
accepting marks give infinitely many accepting visits in the lifted run.
The two inclusions establish language equality.
\end{proof}

\begin{theorem}[Maximal-delay split preservation]
Let $D$ be an input DTA with no pre-existing location invariants. The
maximal-delay variant (Rocq: \texttt{split\_accepts}) and $D$ accept the same timed words.
\end{theorem}
\begin{proof}
Soundness is the projection argument of the previous proof: every copied transition keeps its
label and resets, and its guard is the original guard conjoined with an entry predicate, so
the projection of a run of the split automaton is a run of $D$. For completeness, fix an
accepting run of $D$ and a stay at a noninitial location $q$, entered with valuation $v$
(after the resets of the incoming edge) and left after a delay $\Delta$ by a transition whose guard
clause $c$ holds at the end of the stay. By \texttt{disj\_inv\_cover} there is $k^*$ with
$\operatorname{ub}(c)\Rightarrow U_{k^*}$, so $U_{k^*}$ holds at the end of the stay and
$b-v(x)\ge\Delta$ for every $(x,b)\in U_{k^*}$, that is, $d_{k^*}(v)\ge\Delta$. Let $k$ maximize $d_k(v)$
(\texttt{best\_spec}; ties are all admissible, and $+\infty$ is allowed). Then $d_k(v)\ge d_{k^*}(v)\ge\Delta$,
so every bound of $U_k$ still holds at the end of the stay and, being non-strict upper bounds on increasing
clocks, throughout it; and $G_k(v)$ holds by the choice of $k$, so the incoming edge, whose guard is conjoined with the
reset-adjusted predicate $G_k$, can enter copy $k$. The outgoing transition is available from every
copy. Making this choice at each occurrence lifts the run; accepting marks are inherited, and
the initial occurrence uses the invariant-free initial copy. The two inclusions give language equality.
\end{proof}

The maximal-delay policy establishes a stronger local copy-choice property:
at a given entry valuation, every admitted copy has an admissible residence
duration at least as long as that of any other copy. The helper lemma
\texttt{best\_spec} proves maximization, including $+\infty$ for an empty
conjunction; ties remain admissible. This is neither global deadlock
freedom nor a theorem about time divergence. Compared with all-copy
replication, maximal-delay entry selection can introduce clock-difference
constraints and more DNF clauses, while it may avoid entering copies that
expire strictly earlier. The measured alternatives exhibit a trade-off
(Table~\ref{tab:split-comparison}): on a two-clock example all-copy uses four
more explicit edges but fewer guard atoms and no diagonal constraints; on a
six-branch family both produce 68 edges, while all-copy uses fewer guard
atoms and no diagonal constraints. All-copy is faster on the small case and
slower on the six-branch case in these measurements. These are synthetic
measurements, not a general performance ordering. The default remains the existing
maximal-delay transformation; all-copy is separately verified and exposed
for comparison.

\begin{table*}[t]
\centering
\caption{Measured split policies on direct synthetic TBAs. ``Two-clock'' is the
\texttt{disjunctive\_invariant\_export} record of Table~\ref{tab:synthetic-tba}
and ``six-branch'' the disjunctive-invariant family at $k=6$ of
Table~\ref{tab:parameterized-synthetic}. Locations are graph-reachable
locations. Times are means over 200 in-memory export calls.}
\label{tab:split-comparison}
\scriptsize
\begin{tabular}{@{}llrrrrr@{}}
\toprule
Input & Policy & Reach. loc. & Edges & Guard atoms & Diagonal atoms & Time (ms)\\
\midrule
Two-clock & Maximal delay & 3 & 6 & 12 & 6 & 0.831\\
Two-clock & All-copy & 3 & 10 & 8 & 0 & 0.712\\
Six-branch & Maximal delay & 5 & 68 & 320 & 204 & 15.990\\
Six-branch & All-copy & 5 & 68 & 64 & 0 & 28.731\\
\bottomrule
\end{tabular}
\end{table*}
